% =====================================================================
% Appendix: EmoMCTS qualitative examples and prompts
% Drop-in LaTeX. Requires the preamble block below (once, in your main file).
% =====================================================================

% ---------- PREAMBLE (add once to your main document) ----------------
% \usepackage[most]{tcolorbox}
% \usepackage{listings}
% \usepackage{xcolor}
%
% % Box for self-play dialogues
% \newtcolorbox{dialoguebox}[1][]{
%   colback=gray!4, colframe=gray!55, boxrule=0.4pt, arc=2pt,
%   left=6pt, right=6pt, top=4pt, bottom=4pt, breakable, #1}
% \newcommand{\sys}[1]{\textbf{System} \emph{[#1]}:\ }
% \newcommand{\usr}[1]{\textbf{User} \emph{(#1)}:\ }
% \newcommand{\pivot}{$\rightarrow$\ }   % marks an emotion-driven pivot turn
%
% % Box for prompt templates (verbatim)
% \definecolor{promptbg}{gray}{0.96}
% \lstdefinestyle{prompt}{
%   basicstyle=\ttfamily\scriptsize, breaklines=true, columns=fullflexible,
%   backgroundcolor=\color{promptbg}, frame=single, rulecolor=\color{gray!55},
%   xleftmargin=4pt, xrightmargin=4pt, aboveskip=4pt, belowskip=4pt,
%   commentstyle=\color{gray!70}\itshape, morecomment=[l]{\#}}
% ---------------------------------------------------------------------


% =====================================================================
\section{Qualitative examples}\label{app:examples}

Three self-play dialogues produced by EmoMCTS ($\beta{=}0.7$, top-$K{=}5$, Vicuna-13B). Each
system turn is annotated with the selected dialogue act; each user turn with the classifier
emotion that the parallel channel $Q_{\text{emo}}$ backs up. Turns marked $\rightarrow$ are
emotion-driven pivots, where a negative user emotion prompts a trust-building action that flips
the user back to a positive state. A few system turns are rendered as a short paraphrase of the
chosen act (an open-loop realization artifact), shown faithfully.

\begin{dialoguebox}[title=Example 1: fear about transparency, met with a credibility appeal
(success, 5 turns)]
\small
\sys{greeting} Hello. How are you today?\\[2pt]
\sys{logical appeal} (asks whether the Persuadee has heard of Save the Children)\\
\usr{neutral} I have heard of it, but I'm not familiar with all of the details. Can you tell me more about it?\\[2pt]
\sys{proposition of donation} Save the Children is an organization that helps children in developing countries and war-torn areas. They provide education, healthcare, and protection from harm\ldots\\
\usr{fear} I would like to help, but I am concerned about the transparency and accountability of the organization. I want to make sure my donation will be used effectively.\\[2pt]
\pivot\sys{credibility appeal} I understand your concerns, but I can assure you that Save the Children is a highly reputable organization. They have been around for over 100 years\ldots a strong focus on transparency and accountability, and regularly publish reports on their progress and impact.\\
\usr{happiness} I appreciate the information and the reassurance. I am definitely interested in making a donation to this cause.\\[2pt]
\sys{proposition of donation} (invites the Persuadee to make a small donation)\\
\usr{donate} Yes, I would like to donate \$5 to Save the Children.
\end{dialoguebox}

\noindent\textit{The user's reaction to the proposition is classified as \textbf{fear}.
Instead of repeating the pitch, the planner selects a \textbf{credibility appeal}---the action
whose $Q_{\text{emo}}$ history moves anxious users toward trust---and the user flips to
\textbf{happiness}, then donates.}

\begin{dialoguebox}[title=Example 2: skepticism, empathy, a relapse, and recovery
(success, 7 turns)]
\small
\sys{greeting} Hello. How are you?\\[2pt]
\sys{emotion appeal} (asks about the Persuadee's opinion of Save the Children)\\
\usr{happiness} I'm good, how are you doing?\\[2pt]
\sys{credibility appeal} Very well. I'm just organizing info for my charity called Save the Children. Have you heard of it before?\\
\usr{neutral} No, I have not. Can you tell me more?\\[2pt]
\sys{emotion appeal} Save the Children\ldots helps kids who are in desperate need. They can help with safety, education and more.\\
\usr{surprise} That sounds great. I believe in this charity, but still wonder how much of the money I donate actually helps. I always worry it will just go to some higher-up living the high life.\\[2pt]
\pivot\sys{emotion appeal} Every little bit makes a difference. When people are so poor, it's amazing what a tiny amount can do. I usually donate in hopes I can at least feel like I did my part\ldots\\
\usr{happiness} With that all said, I do feel like any organization that aims to help the children I am more inclined to donate to\ldots they are our future!\\[2pt]
\sys{proposition of donation} I couldn't agree more\ldots if you would like to make a small donation, I would be more than happy to help.\\
\usr{neutral} Hmm, I understand, but I am still not sure. I want to make sure my donation goes towards helping children, not other expenses.\\[2pt]
\pivot\sys{reassurance} I understand your concerns. But I can assure you that Save the Children is a well-established charity\ldots a track record of transparency and accountability\ldots\\
\usr{donate} Alright, I will make a donation of \$5 to Save the Children\ldots
\end{dialoguebox}

\noindent\textit{Two separate dips---\textbf{surprise} (skeptical) and then a \textbf{neutral}
relapse after the proposition---each trigger a non-proposition, trust-building action before the
planner re-attempts the close, showing the emotion channel tracking the user across turns.}

\begin{dialoguebox}[title=Example 3: converting doubt with an emotion appeal (success, 6 turns)]
\small
\sys{greeting} Hey there! How are you today?\\[2pt]
\sys{logical appeal} Save the Children helps children in developing countries\ldots safety, education and more.\\
\usr{surprise} That sounds great. I believe in this charity, but still wonder how much of the money I donate actually helps. I always worry it will just go to some higher up living the high life.\\[2pt]
\pivot\sys{emotion appeal} Every little bit makes a difference\ldots I usually donate in hopes I can at least feel like I did my part. If some corrupt person took it, that's the worst karma and even worse scandal imaginable.\\
\usr{happiness} With that all said I do feel like any organization that aims to help the children I am more inclined to donate to\ldots they are our future!\\[2pt]
\sys{proposition of donation} I completely agree. Children are the future\ldots even a small donation of \$1 can make a difference. Would you like to donate to Save the Children?\\
\usr{donate} Yes, I would like to donate to Save the Children.
\end{dialoguebox}

\noindent\textit{A single \textbf{surprise}/skepticism signal is enough for the planner to
choose an empathetic \textbf{emotion appeal} over a logical argument; the user moves to
\textbf{happiness} and accepts on the next turn.}


% =====================================================================
\section{Prompt templates}\label{app:prompts}

All LLM roles share one Vicuna-13B backbone. Placeholders are in \texttt{\{braces\}};
\texttt{\{in\_context\_example\}} is a single human reference dialogue prepended as a
demonstration (omitted for brevity).

\paragraph{(a) Policy prior --- proposes the next system dialogue act.}
\begin{lstlisting}[style=prompt]
Save the Children is head-quartered in London, and they work to help fight poverty
around the world. Children need help in developing countries and war zones. Small
donations like $1 or $2 go a long way to help.
You are Persuader who is trying to persuade the Persuadee to donate to a charity called
Save the Children. You can choose amongst the following actions during a conversation:
{dialogue_acts}      # [greeting] [task related inquiry] [credibility appeal]
                     # [logical appeal] [emotion appeal] [proposition of donation] ...
{in_context_example}
The following is a new conversation between Persuader (you) and a Persuadee.
{dialogue_history}
Persuader: [
\end{lstlisting}

\paragraph{(b) Value function --- estimates the probability the user will donate.}
\begin{lstlisting}[style=prompt]
The following is background information about the task.
The Persuader is trying to persuade the Persuadee to donate to Save the Children.
The Persuadee can choose amongst the following actions to respond to the Persuader:
{user_dialogue_acts} # [agree donation] [positive reaction] [negative reaction] ...
{in_context_example}
The following is a new conversation between another Persuader and Persuadee.
{dialogue_history}
Persuader: Would you be interested in donating to Save the Children?
Persuadee: [
\end{lstlisting}

\paragraph{(c) User simulator --- produces the persuadee's reaction (the environment).}
\begin{lstlisting}[style=prompt]
You are a persuadee. A Persuader is trying to persuade you to donate to a charity called
Save the Children. You can choose amongst the following actions to respond:
{user_dialogue_acts}
{in_context_example}
The following is a new conversation between a Persuader and a Persuadee (you). You may or
may not want to donate to Save the Children.
{dialogue_history}
Persuadee:
\end{lstlisting}

\paragraph{(d) System utterance (NLG) --- realizes the chosen act into text.}
\begin{lstlisting}[style=prompt]
Save the Children is head-quartered in London ... Small donations like $1 or $2 go a long
way to help. You are Persuader who is trying to persuade the Persuadee to donate to Save
the Children.
{in_context_example}
The following is a new conversation between Persuader (you) and another Persuadee.
{dialogue_history}
Persuader: [{chosen_dialogue_act}]
\end{lstlisting}

\paragraph{(e) Emotion classifier.}
For all reported runs we use the deterministic encoder
\texttt{j-hartmann/emotion-english-distilroberta-base} (no prompt; one forward pass returns the
full distribution over \{anger, disgust, fear, joy$\rightarrow$happiness, neutral, sadness,
surprise\}). A prompt-based alternative is also available:
\begin{lstlisting}[style=prompt]
Classify the emotion in the following utterance into exactly one of:
[Happiness] [Sadness] [Fear] [Anger] [Surprise] [Disgust] [Contempt] [Neutral].
(Each label is given a one-sentence definition.)
Utterance: {user_utterance}
Example Output: [Contempt]
Emotion: [
\end{lstlisting}
